\scanpage{24}
\mspage{2}

\uline{de signe}. Supposons alors $X$ \inserted{lisse,} de dim. relative 1. On peut considérer alors $\underline{\mathrm{Pic}}^d_{X/S}$, \inserted{torseur sous $B=\underline{\mathrm{Pic}}^0_{X/S}$,} correspondant aux faisceaux inv sur $X/S$ de degré relatif $d$. Il y a un morphisme canonique

\[
X\xrightarrow{\psi}\underline{\mathrm{Pic}}^1_{X/S},
\]

\deleted{et celui-ci} défini par la diagonale de $X\times_S X$, qui est un diviseur relatif sur le deuxième facteur.

\diagram{\begin{tikzcd}[ampersand replacement=\&,column sep=large,row sep=large] X\arrow[r,no head,"\operatorname{pr}_1"]\arrow[d,no head]\&X\times_S X\arrow[d,no head,"\operatorname{pr}_2"]\\ S\arrow[r,no head]\&X\end{tikzcd}}

Il fait de $\underline{\mathrm{Pic}}^1_{X/S}$ une solution du pb universel d'Albanese. Il s'ensuit, grâce à la première convention de \unsure{signe|signes}, un isom

\[
c\colon B\xrightarrow{\sim} A=\underline{\mathrm{Pic}}^0_{B/S},
\]

qui s'explicite ainsi : Si $y,x$ sont deux sections de $X$ sur $S$, alors $\psi(y)-\psi(x)$ est une section de $\underline{\mathrm{Pic}}^0_{X/S}$ correspondant au diviseur relatif $y(S)-x(S)$, et on a

\[
c(\psi(y)-\psi(x))=c(\operatorname{cl}(y(S))-\operatorname{cl}(x(S))=N_yN_x^{-1}
\]

où $N$ est un faisceau de Weil sur $X\times_S B$.

\uline{Question à tirer au clair} : Cet isom $c$ définit-il \deleted{\unsure{le}} la $\Theta$-\uline{polarisation} sur $B$? \deleted{\unsure{(ou la symétrique d'une}}\inserted{Non, c'est sans doute la} \deleted{\unsure{polarisation)).}}\inserted{\uline{symétrique}} de la $\Theta$-polarisation.

C'est une question ``géométrique,'' donc on peut supposer $S$ spectre d'un corps alg. clos. On a une polarisation de $B$ par \unsure{le|la} \unsure{classique|d'unique|l'unique} diviseur $\Theta$. L'hom \inserted{correspondant} $u\colon B\longrightarrow A$ est défini par $u(\deleted{\unsure{x}}t)=\text{classe de }t.\Theta-\Theta$. Si
% archive sheet number: 23
